% GENERATED FILE. Edit canonical lessons, then run npm run generate:books.
% canonical-source-hash: fnv1a64:4755f8fe940beff2
% canonical-lessons: HI-C70-game, HI-C70-play, HI-C70-cricket, HI-C70-song, HI-C70-rest

\chapter{After the Work Is Done}
\label{ch:hi-after-the-work-is-done}
\hlchaptermodality{77}

\begin{chapteropening}
\textbf{By the end of this chapter:} \emph{I can say what game I play and how I rest when somebody I have just met asks what I do.}

Complete HI-C70-rest and say all five words in order, then say what game you play.
\end{chapteropening}

\section[khel]{\hi{खेल} (khel) — a game}

\label{lesson:HI-C70-game}



\begin{quote}
Two words back: what does \emph{ṭikaṭ} mean, and what is \textbf{\hi{जाना}}? The next five words are what you do when you are not going anywhere.
\end{quote}

\subsection*{You'll want to know: \hi{खेल}}
\textbf{\hi{खेल}} (\emph{khel}) — \textquotedblleft{}a game\textquotedblright{}.

Sanskrit \hi{खेल} (\emph{khela}), \textquotedblleft{}play, sport\textquotedblright{}, from a root that meant shaking or moving about. The final vowel dropped away by the schwa rule and left the bare consonant ending Hindi likes.

\hi{खेल} covers more than a sport with rules. A child's game is a \hi{खेल}, a match is a \hi{खेल}, and a thing somebody is not taking seriously is a \hi{खेल} too. When somebody you have just met asks what you do with your evenings, this is the word the answer starts from.

The \textbf{\hi{ख}} is the aspirated one — a real puff of breath after the \emph{k}, which is what separates \hi{खेल} from \hi{क}, and the difference changes the word.

\begin{grammarlens}[title={What you've built}]
The first of five words for the hours after work.
\end{grammarlens}

\subsection*{Your turn}
\begin{itemize}
\raggedright
  \item \emph{Recall:} say \emph{ṭikaṭ}, then read \textbf{\hi{जाना}} and say what it means
  \item \emph{Hear:} \emph{khel}, with the breath on the first sound
  \item \emph{Say it:} \emph{khel}
  \item \emph{Contrast:} \emph{khel} said with the puff, and the same shape said without it
  \item \emph{Recall:} say \emph{desh}, then read \textbf{\hi{गाड़ी}}
\end{itemize}

\begin{tcolorbox}[breakable,colback=teal!4,colframe=teal!35!black,title={Before you move on}]
What does \hi{खेल} mean? (\textquotedblleft{}a game\textquotedblright{}.) Say it once more.
\end{tcolorbox}
\section[khelnā]{\hi{खेलना} (khelnā) — to play}

\label{lesson:HI-C70-play}



\begin{quote}
Say \emph{khel}, then write \textbf{\hi{खेल}}: \hi{ख}, then \hi{े} above it, then \hi{ल}.
\end{quote}

\subsection*{You'll want to know: \hi{खेलना}}
\textbf{\hi{खेलना}} (\emph{khelnā}) — \textquotedblleft{}to play\textquotedblright{}.

The same \hi{खेल} you learnt in the last lesson, with \textbf{-\hi{ना}} on the end. That is the whole of it: Hindi makes a doing word out of a naming word by adding those two letters, and you have now watched it happen once.

Look at what that gives you. \textbf{\hi{जाना}} was built the same way, and so were \textbf{\hi{बोलना}}, \textbf{\hi{लिखना}}, \textbf{\hi{सुनना}}. Four you already own and one you have just made — the ending is a tool, not a list to memorise.

\textbf{\hi{मैं} \hi{खेलता} \hi{हूँ}} is \textquotedblleft{}I play\textquotedblright{}, in the frame you have used since \textbf{\hi{मैं} \hi{हिंदी} \hi{बोलता} \hi{हूँ}}.

\begin{grammarlens}[title={What you've built}]
A game, and the doing word that comes out of it.
\end{grammarlens}

\subsection*{Your turn}
\begin{itemize}
\raggedright
  \item \emph{Recall:} read \textbf{\hi{जाना}}, then say it without looking
  \item \emph{Say it:} \emph{khel}, then \emph{khelnā}
  \item \emph{Say it:} \emph{maiṁ kheltā hūṁ}
  \item \emph{Write it:} \hi{खेल}, then add \hi{ना} to the end of it
  \item \emph{Recall:} read \textbf{\hi{भारत}}, then say \emph{rel}
\end{itemize}

\begin{tcolorbox}[breakable,colback=teal!4,colframe=teal!35!black,title={Before you move on}]
What does \hi{खेलना} mean? (\textquotedblleft{}to play\textquotedblright{}.) And \hi{खेल}? (\textquotedblleft{}a game\textquotedblright{}.)
\end{tcolorbox}
\section[krikeṭ]{\hi{क्रिकेट} (krikeṭ) — cricket, the game}

\label{lesson:HI-C70-cricket}



\begin{quote}
Which of your two words is the game, and which is the doing word made out of it?
\end{quote}

\subsection*{You'll want to know: \hi{क्रिकेट}}
\textbf{\hi{क्रिकेट}} (\emph{krikeṭ}) — \textquotedblleft{}cricket\textquotedblright{}.

English, borrowed whole, and written with two things you have already met. The opening \hi{क्र} is \hi{क} stacked on \hi{र} by the same vowel-killer \hi{्} that stacks \hi{स्ट} in \hi{स्टेशन}. The final \hi{ट} is retroflex, because an English \emph{t} always lands there.

The \hi{ि} after the stack is the short \emph{i} matra, and it sits to the \textbf{left} of the whole cluster even though it is sounded after it — the same preposed matra you first met in a name.

If you learn one game word in Hindi, learn this one. \textbf{\hi{मैं} \hi{क्रिकेट} \hi{खेलता} \hi{हूँ}} puts three of this chapter's ideas into a single sentence.

\begin{grammarlens}[title={What you've built}]
Three: \hi{खेल}, \hi{खेलना}, \hi{क्रिकेट}. A game, playing it, and one to name.
\end{grammarlens}

\subsection*{Your turn}
\begin{itemize}
\raggedright
  \item \emph{Read it:} \hi{क्रिकेट}, and find the \hi{ि} sitting to the left of the stack
  \item \emph{Say it:} \emph{krikeṭ}
  \item \emph{Say it:} \emph{maiṁ krikeṭ kheltā hūṁ}
  \item \emph{Recall:} say \emph{shahar}, then read \textbf{\hi{स्टेशन}}
\end{itemize}

\begin{tcolorbox}[breakable,colback=teal!4,colframe=teal!35!black,title={Before you move on}]
What does \hi{क्रिकेट} mean? (\textquotedblleft{}cricket\textquotedblright{}.) Say that you play it.
\end{tcolorbox}
\section[gānā]{\hi{गाना} (gānā) — a song}

\label{lesson:HI-C70-song}



\begin{quote}
Write \textbf{\hi{क्रिकेट}} from memory, starting with the \hi{ि} you put down first.
\end{quote}

\subsection*{You'll want to know: \hi{गाना}}
\textbf{\hi{गाना}} (\emph{gānā}) — \textquotedblleft{}a song\textquotedblright{}.

Sanskrit \hi{गान} (\emph{gāna}), \textquotedblleft{}singing\textquotedblright{}, from \hi{गै} (\emph{gai-}), \textquotedblleft{}to sing\textquotedblright{}. The root is one of the oldest in the language and it has never stopped working.

Now notice the shape. \hi{गाना} ends in \textbf{-\hi{ना}}, exactly like \hi{खेलना} — and that is not an accident. \hi{गाना} is also the doing word \textquotedblleft{}to sing\textquotedblright{}. One spelling, two jobs: the song, and the singing of it. Which one is meant is settled by the sentence around it and never by the word.

Hindi does this often enough that it is worth expecting rather than being surprised by. You will meet another word with the same double life before this run of chapters is out.

\begin{grammarlens}[title={What you've built}]
Four, and the first word you have met that is two words at once.
\end{grammarlens}

\subsection*{Your turn}
\begin{itemize}
\raggedright
  \item \emph{Say it:} \emph{gānā}
  \item \emph{Contrast:} \emph{gānā} as the song, and \emph{gānā} as the singing
  \item \emph{Write it:} \hi{गाना}, then \hi{खेलना} under it, and look at the ending they share
  \item \emph{Recall:} read \textbf{\hi{भाषा}}, then say \emph{ṭikaṭ}
\end{itemize}

\begin{tcolorbox}[breakable,colback=teal!4,colframe=teal!35!black,title={Before you move on}]
What two things can \hi{गाना} mean? (\textquotedblleft{}a song\textquotedblright{}, and \textquotedblleft{}to sing\textquotedblright{}.)
\end{tcolorbox}
\section[ārām]{\hi{आराम} (ārām) — rest, ease}

\label{lesson:HI-C70-rest}



\begin{quote}
Name the four words of this chapter so far.
\end{quote}

\subsection*{You'll want to know: \hi{आराम}}
\textbf{\hi{आराम}} (\emph{ārām}) — \textquotedblleft{}rest, ease\textquotedblright{}.

Persian \ar{آرام} (\emph{ārām}), \textquotedblleft{}quiet, calm, repose\textquotedblright{}. It came in with the Persian half of the vocabulary and settled deep: \textbf{\hi{आरामकुर्सी}} is an easy chair, built from \hi{आराम} and the \textbf{\hi{कुर्सी}} you already own.

\hi{आराम} is not sleep. \textbf{\hi{नींद}} is sleep, and you have that word. \hi{आराम} is the state of not being pushed — sitting down, breathing out, nothing owed to anybody for an hour. \textbf{\hi{आराम} \hi{करना}} is what you say for taking that hour.

It is also an answer. When somebody asks \textbf{\hi{आप} \hi{कैसे} \hi{हैं}?}, one ordinary reply is that you are at \hi{आराम}, and it says more warmly than \textbf{\hi{ठीक}} does.

\begin{grammarlens}[title={What you've built}]
Five: \hi{खेल}, \hi{खेलना}, \hi{क्रिकेट}, \hi{गाना}, \hi{आराम}. Enough to say what you do when the day lets go of you.
\end{grammarlens}

\subsection*{Your turn}
\begin{itemize}
\raggedright
  \item \emph{Say it:} \emph{ārām}
  \item \emph{Say it:} \emph{ārām karnā}
  \item \emph{Write it:} \hi{गाना} and then \hi{आराम}, both from memory
  \item \emph{Say it:} all five in order
  \item \emph{Recall:} say \emph{angrezī}, then read \textbf{\hi{जाना}}
\end{itemize}

\begin{tcolorbox}[breakable,colback=teal!4,colframe=teal!35!black,title={Before you move on}]
What does \hi{आराम} mean? (\textquotedblleft{}rest\textquotedblright{}.) And \hi{गाना}? (\textquotedblleft{}a song\textquotedblright{}, or \textquotedblleft{}to sing\textquotedblright{}.)
\end{tcolorbox}
